\documentclass[11pt]{article}

% ===================== PACKAGES =====================
\usepackage[a4paper,margin=1in]{geometry}
\usepackage{amsmath,amssymb}
\usepackage{enumitem}
\usepackage{fancyhdr}
\usepackage{xcolor}
\usepackage{setspace}
\usepackage{hyperref}

% ===================== HEADER & FOOTER =====================
\pagestyle{fancy}
\fancyhf{}
\lhead{CSE 400: Fundamentals of Probability in Computing}
\rhead{Lecture 10: Min-Cut Algorithms}
\cfoot{\thepage}

% ===================== DOCUMENT =====================
\begin{document}

\begin{center}
    {\Large \textbf{CSE 400 – Lecture Scribe}}\\[0.2cm]
    {\normalsize School of Engineering and Applied Sciences, Ahmedabad University}\\[0.3cm]
    {\large \textbf{Lecture 10: Min-Cut Algorithms (Deterministic and Randomized)}}
\end{center}

\vspace{0.4cm}

\section*{1. Min-Cut Problem}

\subsection*{1.1 Why Use Min-Cut?}

\textbf{Step-by-step reasoning (as per lecture):}
\begin{enumerate}
    \item The min-cut algorithm is used to solve problems related to:
    \begin{itemize}
        \item Network connectivity
        \item Reliability
        \item Optimization
    \end{itemize}
    \item In \textbf{network design}, min-cut:
    \begin{itemize}
        \item Helps improve efficiency of communication
        \item Optimizes network flow
        \item Finds the minimum capacity cut
    \end{itemize}
    \item In \textbf{communication networks}, min-cut:
    \begin{itemize}
        \item Helps understand vulnerability of networks to failures
        \item Assists in building robust and fault-tolerant networks
    \end{itemize}
    \item In \textbf{VLSI design}, min-cut:
    \begin{itemize}
        \item Helps partition circuits into smaller components
        \item Reduces interconnectivity complexity
    \end{itemize}
\end{enumerate}

\subsection*{1.2 What Is Min-Cut?}

\textbf{Definitions (step-by-step):}
\begin{enumerate}
    \item A \textbf{cut-set} in a graph is:
    \begin{itemize}
        \item A set of edges whose removal breaks the graph into two or more connected components
    \end{itemize}
    \item Given a graph
    \[
    G = (V, E)
    \]
    with \( n \) vertices:
    \begin{itemize}
        \item The \textbf{minimum cut (min-cut)} problem is to find a cut-set with \textbf{minimum cardinality}
    \end{itemize}
    \item Min-cut algorithms (e.g., Karger’s algorithm):
    \begin{itemize}
        \item Are random
        \item Can be sensitive to the initial choice of edges
    \end{itemize}
    \item If critical edges are contracted early:
    \begin{itemize}
        \item The algorithm may find a smaller cut
    \end{itemize}
\end{enumerate}

\subsection*{1.3 Edge Contraction (Core Operation)}

\textbf{Step-by-step operation:}
\begin{enumerate}
    \item The main operation is \textbf{edge contraction}
    \item Contracting an edge \( (u, v) \):
    \begin{itemize}
        \item Merge vertices \( u \) and \( v \) into a single vertex
        \item Eliminate all edges connecting \( u \) and \( v \)
        \item Retain all other edges
    \end{itemize}
    \item After contraction:
    \begin{itemize}
        \item Parallel edges may exist
        \item Self-loops do not exist
    \end{itemize}
\end{enumerate}

\subsection*{1.4 Successful Min-Cut Run}

\textbf{Step-by-step definition:}
\begin{enumerate}
    \item A successful min-cut run:
    \begin{itemize}
        \item Refers to success in the outcome of an algorithm
    \end{itemize}
    \item The algorithm:
    \begin{itemize}
        \item Correctly identifies the minimum cut of the graph
    \end{itemize}
    \item The resulting cut corresponds to the true minimum cut
\end{enumerate}

\subsection*{1.5 Unsuccessful Min-Cut Run}

\textbf{Step-by-step definition:}
\begin{enumerate}
    \item An unsuccessful min-cut run:
    \begin{itemize}
        \item Refers to an iteration of a min-cut algorithm
    \end{itemize}
    \item The algorithm:
    \begin{itemize}
        \item Fails to correctly identify the minimum cut
    \end{itemize}
    \item This occurs due to:
    \begin{itemize}
        \item Unfavorable contraction choices during execution
    \end{itemize}
\end{enumerate}

\subsection*{1.6 Max-Flow Min-Cut Theorem}

\textbf{Theorem statement (as given):}

\begin{quote}
“In a flow network, the maximum amount of flow passing from the source to the sink is equal to the total weight of the edges in a minimum cut.”
\end{quote}

\textbf{Step-by-step definitions:}
\begin{enumerate}
    \item \textbf{Capacity of a cut}:
    \begin{itemize}
        \item Sum of capacities of edges oriented from vertex \( \in X \) to vertex \( \in Y \)
    \end{itemize}
    \item \textbf{Minimum cut}:
    \begin{itemize}
        \item The cut with the smallest possible capacity
    \end{itemize}
    \item \textbf{Minimum cut capacity}:
    \begin{itemize}
        \item Capacity of the minimum cut
    \end{itemize}
    \item \textbf{Maximum flow}:
    \begin{itemize}
        \item Largest possible flow from source \( S \) to sink \( T \)
    \end{itemize}
\end{enumerate}

\section*{2. Deterministic Min-Cut Algorithm}

\subsection*{2.1 Stoer–Wagner Minimum Cut Algorithm}

\textbf{Theorem-based reasoning (step-by-step):}
\begin{enumerate}
    \item Let \( s \) and \( t \) be two vertices of graph \( G \)
    \item Let \( G/\{s,t\} \) be the graph obtained by merging \( s \) and \( t \)
    \item A minimum cut of \( G \) is:
    \begin{itemize}
        \item The smaller of:
        \begin{itemize}
            \item A minimum \( s\text{-}t \) cut of \( G \)
            \item A minimum cut of \( G/\{s,t\} \)
        \end{itemize}
    \end{itemize}
    \item Two cases:
    \begin{itemize}
        \item Case 1: A minimum cut separates \( s \) and \( t \)
        \begin{itemize}
            \item Then minimum \( s\text{-}t \) cut is a minimum cut
        \end{itemize}
        \item Case 2: No minimum cut separates \( s \) and \( t \)
        \begin{itemize}
            \item Then minimum cut of \( G/\{s,t\} \) is the answer
        \end{itemize}
    \end{itemize}
\end{enumerate}

\subsection*{2.2 Deterministic Min-Cut Algorithm – Pseudocode}

\textbf{Algorithm 1: \texttt{MinimumCutPhase(G, a)}}
\begin{verbatim}
A ← {a}
while A \neq V do
    add to A the most tightly connected vertex
return the cut weight as the “cut of the phase”
\end{verbatim} 

\textbf{Algorithm 2: \texttt{MinimumCut(G)}}
\begin{verbatim}
while |V| \geq 1 do
    choose any a from V
    MinimumCutPhase(G, a)
    if the cut-of-the-phase is lighter than the current minimum cut then
        store the cut-of-the-phase as the current minimum cut
    shrink G by merging the two vertices added last
return the minimum cut
\end{verbatim} 

\section*{3. Randomized Min-Cut Algorithm}

\subsection*{3.1 Why Randomized Algorithm?}

\textbf{Step-by-step points from lecture:}
\begin{enumerate}
    \item Randomized algorithms:
    \begin{itemize}
        \item Provide probabilistic guarantee of success
    \end{itemize}
    \item They:
    \begin{itemize}
        \item Provide a more accurate estimate with fewer iterations
    \end{itemize}
    \item Additional listed properties:
    \begin{itemize}
        \item Efficiency
        \item Parallelization
        \item Approximation guarantees
        \item Avoidance of worst-case instances
        \item Heuristic nature
        \item Robustness
    \end{itemize}
\end{enumerate}

\subsection*{3.2 Karger’s Randomized Algorithm}

\textbf{Algorithmic idea (as shown):}
\begin{enumerate}
    \item Randomly pick an edge
    \item Remove it
    \item Fuse its two endpoints
    \item Repeat contraction until only two vertices remain
    \item The remaining edges form a cut
    \item The output cut may or may not be the minimum cut
\end{enumerate}

\subsection*{3.3 Randomized Min-Cut Algorithm – Pseudocode}

\textbf{Algorithm 3: \texttt{RECURSIVE-RANDOMIZED-MIN-CUT(G, \(\alpha\))}}
\begin{verbatim}
Input: undirected multigraph G with n vertices, integer alpha > 0
Output: a cut C of G
if n <= alpha^3 then
    C <- a min-cut of G found using brute force (exhaustive search)
else
    for i <- 1 to alpha do
        G' <- multigraph obtained by applying n - ceil(n / sqrt(alpha)) random contraction steps on G
        C' <- RECURSIVE-RANDOMIZED-MIN-CUT(G', alpha)
        if i = 1 or |C'| < |C| then
            C <- C'
return C
\end{verbatim}

\subsection*{3.4 Comparison: Deterministic vs Randomized Min-Cut Algorithm}

\textbf{Decision basis:}
\begin{itemize}
    \item Any specific problem is responsible for deciding the approach
\end{itemize}

\textbf{Deterministic Min-Cut:}
\begin{enumerate}
    \item Always guarantees exact minimum cut
    \item May have higher time complexity for large graphs
    \item Stoer–Wagner time complexity:
    \[
    O(V \cdot E + V^2 \log V)
    \]
\end{enumerate}

\textbf{Randomized Min-Cut:}
\begin{enumerate}
    \item Provides approximate minimum cut with high probability
    \item Karger’s algorithm time complexity:
    \[
    O(V^2)
    \]
\end{enumerate}

\subsection*{3.5 Theorem for Min-Cut Set}

\textbf{Theorem (as stated):}
\[
\text{The algorithm outputs a min-cut set with probability at least } \frac{2}{n(n-1)}
\]

\subsection*{3.6 Python Simulation}

\textbf{Lecture instruction:}
\begin{enumerate}
    \item Open Campuswire post regarding Lecture 10
    \item Download the provided \texttt{.ipynb} file
    \item Use it for simulation of the randomized min-cut algorithm
\end{enumerate}

\subsection*{3.7 Class Participation}

\begin{itemize}
    \item Participation linked to:
    \begin{itemize}
        \item Campuswire discussion
        \item Python simulation activity
    \end{itemize}
\end{itemize}

\vspace{0.3cm}
\begin{center}
\textbf{End of Lecture Scribe}
\end{center}

\end{document}
